% Part: incompleteness
% Chapter: representability-in-q
% Section: introduction

\documentclass[../../../include/open-logic-section]{subfiles}

\begin{document}

\olfileid{inc}{req}{int}
\olsection{परिचय}

अपूर्णता प्रमेये अशा उपपत्तींना लागू होतात ज्यांत संगणनक्षम
फलनांविषयीची मूलभूत तथ्ये व्यक्त आणि सिद्ध करता येतात. अशा प्रकारची
अगदी किमान उपपत्ती आपण पाहणार आहोत; तिला ``$\Th{Q}$'' म्हणतात
(राफेल रॉबिन्सन यांच्या नावावरून तिला कधीकधी ``रॉबिन्सनची $Q$''
असेही म्हणतात). एखादे फलन $\Th{Q}$ मध्ये \emph{निरूपणीय} असणे म्हणजे
काय, हे आपण सांगू आणि मग पुढील विधान सिद्ध करू:
\begin{quote}
  फलन $\Th{Q}$ मध्ये निरूपणीय असते, तर आणि तरच ते संगणनक्षम असते.
\end{quote}
यामुळे आपल्याला संगणनक्षमतेचे आणखी एक प्रतिमान मिळते. शिवाय,
थांबण्याच्या समस्येचे त्याकडे न्यूनीकरण करून
$\Setabs{!A}{\Th{Q} \Proves !A}$ हा संच निर्णेय नाही, हे दाखवण्यासाठीही
त्याचा उपयोग करू. पुढे याहून अधिक बलवान निष्कर्ष आपण सिद्ध करू.

$\Th{Q}$ ची भाषा ही अंकगणिताची भाषा आहे; $\Th{Q}$ मध्ये पुढील
स्वयंसिद्धके आहेत. ती !!{identity} असलेल्या प्रथम-क्रम तर्कशास्त्रातील
इतर स्वयंसिद्धके आणि नियम यांच्यासोबत वापरायची आहेत:
\begin{align*}
& \lforall[x][\lforall[y][(\eq[x'][y'] \lif \eq[x][y])]] \tag{$!Q_1$}\\
& \lforall[x][\eq/[\Obj 0][x']] \tag{$!Q_2$}\\
& \lforall[x][(\eq[x][\Obj 0] \lor \lexists[y][\eq[x][y']])] \tag{$!Q_3$}\\
& \lforall[x][\eq[(x + \Obj 0)][x]] \tag{$!Q_4$}\\
& \lforall[x][\lforall[y][\eq[(x + y')][(x + y)']]] \tag{$!Q_5$}\\
& \lforall[x][\eq[(x \times \Obj 0)][\Obj 0]] \tag{$!Q_6$}\\
& \lforall[x][\lforall[y][\eq[(x \times y')][((x \times y) + x)]]] \tag{$!Q_7$}\\
& \lforall[x][\lforall[y][(x < y \liff \lexists[z][\eq[(z' + x)][y]])]] \tag{$!Q_8$}
\end{align*}
प्रत्येक नैसर्गिक संख्या $n$ साठी $\num{n}$ हा संख्यांक
$\Obj{0}^{\prime\prime\ldots\prime}$ हे पद म्हणून परिभाषित करा;
त्यात एकूण $n$ उत्तरवर्ती चिन्हे असतात. म्हणजे $\num{0}$ हा
!!{constant}~$\Obj{0}$ एकटाच, $\num{1}$ हे $\Obj{0}'$,
$\num{2}$ हे $\Obj{0}''$, इत्यादी.

अंकगणिताची उपपत्ती म्हणून $\Th{Q}$ \emph{अत्यंत} दुर्बल आहे;
उदाहरणार्थ, $\lforall[x][\eq/[x][x']]$ किंवा
$\lforall[x][\lforall[y][(x + y) = (y + x)]]$ यांसारखी
अगदी साधी तथ्येसुद्धा त्यात सिद्ध करता येत नाहीत. पण $\Th{Q}$
इतकी रंजक असण्याचे मोठे कारण ती \emph{दुर्बल} आहे, हेच आहे.
अपूर्णता प्रमेय लागू होण्याइतकीच तिची शक्ती आहे. $\Th{Q}$
रंजक असण्याचे आणखी एक कारण म्हणजे तिचा स्वयंसिद्धकांचा संच
\emph{सांत} आहे.

$\Th{Q}$ पेक्षा अधिक बलवान उपपत्ती \emph{पेआनो अंकगणित}
$\Th{PA}$ हिला $\Th{Q}$ मध्ये विगमन-रूपबंध जोडून मिळवतात:
\[
(!A(\Obj 0) \land \lforall[x][(!A(x) \lif !A(x'))]) \lif \lforall[x][!A(x)]
\]
येथे $!A(x)$ हे कोणतेही सूत्र असते. $!A(x)$ मध्ये $x$ व्यतिरिक्त
इतर !!{variable}s मुक्त असतील, तर त्या सर्वांना बांधण्यासाठी
आरंभी सार्वत्रिक संख्यक जोडतो (म्हणजे विगमन-रूपबंधाचे संबंधित
उदाहरण !!a{sentence} असेल). उदाहरणार्थ, $!A(x, y)$ मध्ये
!!{variable}~$y$ हेसुद्धा मुक्त असेल, तर संबंधित उदाहरण असे आहे:
\[
\lforall[y][((!A(\Obj 0) \land \lforall[x][(!A(x) \lif !A(x'))]) \lif
  \lforall[x][!A(x)])]
\]
विगमन-रूपबंधाची उदाहरणे वापरून $\Th{Q}$ च्या स्वयंसिद्धकांपेक्षा
$\Th{PA}$ च्या स्वयंसिद्धकांपासून बरेच अधिक सिद्ध करता येते.
खरे तर, पेआनो अंकगणितात सिद्ध करता येत नाहीत अशी नैसर्गिक
संख्यांविषयीची ``स्वाभाविक'' विधाने शोधणे बरेच कठीण असते!{}

\begin{defn}
\ollabel{defn:representable-fn}
  नैसर्गिक संख्यांपासून नैसर्गिक संख्यांकडे जाणारे
  $f(x_0,\ldots,x_k)$ हे फलन $\Th{Q}$ मध्ये
  {\em निरूपणीय} आहे असे म्हणतात, जर असे सूत्र
  $!A_f(x_0,\dots,x_k,y)$ असेल की
  $f(n_0,\dots,n_k) = m$ असेल तेव्हा $\Th{Q}$ पुढील दोन्ही
  विधाने सिद्ध करते:
\begin{enumerate}
\item\ollabel{defn:rep:a} $!A_f(\num{n_0}, \dots, \num{n_k}, \num{m})$
\item\ollabel{defn:rep:b} $\lforall[y][(!A_f(\num{n_0}, \dots,
\num{n_k}, y) \lif \num{m} = y)]$.
\end{enumerate}
\end{defn}

ही व्याख्या इतर मार्गांनीही मांडता येते. उदाहरणार्थ, $\Th{Q}$
$\lforall[y][(!A_f(\num{n_0}, \dots,
\num{n_k}, y) \liff \eq[y][\num{m}])]$ सिद्ध करते, अशी
समतुल्य अटही आपण ठेवू शकलो असतो.

\begin{thm}
\ollabel{thm:representable-iff-comp}
फलन $\Th{Q}$ मध्ये निरूपणीय असते, तर आणि तरच ते संगणनक्षम असते.
\end{thm}

या प्रमेयाच्या पुराव्याला दोन दिशा आहेत. विन्यासमीमांसेचे अंकगणितीकरण
झाल्यावर डावीकडून उजवीकडे जाणारी दिशा तुलनेने सोपी आहे. दुसऱ्या
दिशेसाठी अधिक काम लागते. मूलभूत कल्पना अशी: ``संगणनक्षम'' याचा
अर्थ नेमका करण्यासाठी ``सामान्य पुनरावर्ती'' हा पर्याय निवडतो आणि
प्रत्येक सामान्य पुनरावर्ती फलन $\Th{Q}$ मध्ये निरूपणीय आहे,
हे दाखवतो. लक्षात घ्या, शून्य फलन~$\Zero$, उत्तरवर्ती फलन~$\Succ$
आणि प्रक्षेपण फलने~$\Proj{n}{i}$ यांच्यापासून संयोजन, आदिम पुनरावर्तन
आणि नियमित लघुतमीकरण वापरून परिभाषित करता येणारे फलन म्हणजे
सामान्य पुनरावर्ती फलन. म्हणून प्रत्येक सामान्य पुनरावर्ती फलन
$\Th{Q}$ मध्ये निरूपणीय आहे हे दाखवण्याचा एक मार्ग असा: मूलभूत
फलने निरूपणीय आहेत हे दाखवायचे आणि काही फलने निरूपणीय असताना
त्यांच्यापासून संयोजन, आदिम पुनरावर्तन किंवा नियमित लघुतमीकरणाने
बनणारी फलनेही निरूपणीय असतात, हे दाखवायचे. दुसऱ्या शब्दांत,
मूलभूत फलने निरूपणीय आहेत आणि निरूपणीय फलनांचा वर्ग संयोजन,
आदिम पुनरावर्तन व नियमित लघुतमीकरण यांखाली ``बंद'' आहे,
हे दाखवता येईल. त्यावरून प्रत्येक सामान्य पुनरावर्ती फलन
निरूपणीय असल्याची हमी मिळते.

पण निरूपणीय फलनांचा वर्ग आदिम पुनरावर्तनाखाली बंद आहे हे दाखवण्याची
पायरी कठीण ठरते. ती टाळण्यासाठी, आधी आपण दाखवतो की प्रत्यक्षात
आदिम पुनरावर्तनाशिवायही काम चालते. म्हणजे, प्रत्येक सामान्य
पुनरावर्ती फलन मूलभूत फलनांपासून केवळ संयोजन आणि नियमित
लघुतमीकरण वापरून परिभाषित करता येते, हे दाखवतो. त्यासाठी आदिम
पुनरावर्तन एका विशिष्ट नियमित लघुतमीकरणाने करता येते, हे दाखवतो.
मात्र हे साधण्यासाठी काही आणखी मूलभूत फलने घ्यावी लागतात:
बेरीज, गुणाकार आणि एकरूपता संबंधाचे जातिबोधक फल~$\Char{=}$. मग
या \emph{सर्व} मूलभूत फलनांची $\Th{Q}$ मधील निरूपणीयता आणि
निरूपणीय फलनांचा संयोजन व नियमित लघुतमीकरण यांखालील बंदपणा
दाखवून हे प्रमेय सिद्ध करता येते.

\end{document}
